% !PRINTED 738
\weberpage{738}

For a better overview we set down the equations (13) here once
more without the index $k$, which distinguishes the conjugate values
from one another:
\weq{14}{
\begin{aligned}
2\, l_i & = m_{1,i}\, l'_1 + m_{2,i}\, l'_2 + \cdots + m_{\nu'-1,i}\, l'_{\nu'-1} \\
& \quad + M_{0,i}\, L_0 + M_{1,i}\, L_1 + M_{2,i}\, L_2 + \cdots + M_{\nu'-1,i}\, L_{\nu'-1}.
\end{aligned}
}

If we now form the regulator $E$ according to (13), and denote
by $M$ the absolute value of the determinant of the $\nu-1$ rows
of the numbers $m_{k,i}$, $M_{k,i}$:
\[
M = \pm \Sigma \pm m_{1,1} \ldots m_{\nu'-1,\nu'-1}\, M_{0,0}\, M_{1,1} \ldots M_{\nu'-1,\nu'-1},
\]
then it follows:
\weq{15}{2^{\nu-1} E = M R,}
and from this, having regard to (12):
\weq{16}{E = 2^{-\nu'} M E' T_0.}

It can now be shown further that $M$ is a power of 2 and is
divisible by $2^{\nu'}$.

For since the system (1) was assumed to be a fundamental system
in $H_m$, there exist \gesperrt{integral rational numbers}
$n_{k,i}$, $N_{k,i}$ which satisfy the equations:
\weq{17}{
\begin{aligned}
l'_i & = n_{1,i}\, l_1 + n_{2,i}\, l_2 + \cdots + n_{\nu-1,i}\, l_{\nu-1}, \quad i = 1, 2, \ldots, \nu'-1 \\
L_i & = N_{1,i}\, l_1 + N_{2,i}\, l_2 + \cdots + N_{\nu-1,i}\, l_{\nu-1}, \quad i = 0, 1, 2, \ldots, \nu'-1.
\end{aligned}
}

If we form from these the determinant $R$ and denote by $N$
the absolute value of the determinant of the $n_{k,i}$, $N_{k,i}$, then it follows:
\[
R = N E,
\]
and from this, by (15):
\[
M N = 2^{\nu-1};
\]
from this it follows that each of the two integers $M$, $N$ must
be a power of 2.

In order to decide by which power of 2 the determinant
$M$ is divisible, we first determine a system of integral
rational numbers $a_i$ without a common divisor, which satisfies
the linear equations
\weq{18}{\sum\limits_{1,\, \nu-1}^{i} a_i\, m_{k,i} = 0 \quad (k = 1, 2, \ldots, \nu'-1).}
If we then put
\[
\sum\limits_{1,\, \nu-1}^{i} a_i\, M_{k,i} = \xi_k,
\]
%CONT>>
